

Motion forecasting is central to visual intelligence: agents must anticipate how objects will move in order to plan actions, reason about physical interactions, and synthesize realistic futures. We argue that 3D points in world coordinates provide a general representation that is class-agnostic, view-stable, compact, and directly useful for downstream tasks.
We formalize the task of \textit{goal-conditioned 3D point motion forecasting}: given a short visual history, a set of 3D query points on an object of interest, and a language description of the intended goal, the model predicts the future 3D trajectory of each point. We introduce a full stack to study this task at scale: (1) \textbf{MolmoMotion-1M} is a large corpus of action-described, object-grounded 3D point trajectory dataset annotated from 1.16M unconstrained videos; (2) \textbf{\benchmarkname{}} is a human-verified benchmark spanning 111 object categories and 61 motion types; and (3) \textbf{\modelname} is a general motion forecasting model that supports both autoregressive coordinate prediction and flow-matching-based trajectory generation. 
\modelname is able to accurately predicts diverse motion patterns with different language instructions, and significantly outperforms all existing motion prediction baselines on \benchmarkname{}. Finally, we show that the learned 3D motion prior transfers well to downstream applications: it improves training efficiency and generalization for robot manipulation, 
and its predicted trajectories provide effective motion guidance for generative models to synthesize videos with more realistic object motion.
